\documentclass[12pt, a4paper]{article}
\usepackage[utf8]{utf8}
\usepackage[russian]{babel}
\usepackage{amsmath, amsfonts, amssymb}
\usepackage[margin=2cm]{geometry}

\begin{document}

\section*{Вариант 1}

\begin{enumerate}
    \item $f(x,y) = x + y, \quad x^2 + y^2 = 2 \quad \text{найти } \max \text{ и } \min$
    
    Ограничение: $h(x,y) = x^2 + y^2 - 2 = 0$

    \item $f(x,y) = x + y$ и $h(x,y) = x^2 + y^2 - 2$ определены и непрерывно дифференцируемы на всей $\mathbb{R}^2$.
    
    Пусть $\quad x = \sqrt{2}, \quad y = 0 \Rightarrow (\sqrt{2})^2 + 0^2 = 2 \Rightarrow$ существует хотя бы одно действительное решение.

    \subsection*{Проверка регулярности:}
    $h(x,y) = x^2 + y^2 - 2 = 0$
    
    $\dfrac{\partial h}{\partial x} = 2x, \quad \dfrac{\partial h}{\partial y} = 2y$
    
    $\nabla h(x,y) = \begin{pmatrix} 2x \\ 2y \end{pmatrix} \Rightarrow$ градиент не равен 0 во всех допустимых точках.
    
    Если бы он был нулевым, то $2x = 0$ и $2y = 0 \Rightarrow$ не удовлетворяет $x^2 + y^2 = 2$.
    
    $\Rightarrow$ \textbf{ограничение регулярно}.

    \vspace{0.5cm}
    $D = \{(x,y) \in \mathbb{R}^2 : x^2 + y^2 = 2\}$ — окружность $\Rightarrow$ множество замкнутое и ограниченное $\Rightarrow$ \textbf{компактное}.
    
    Целевая функция непрерывна на всем $\mathbb{R}^2$, а значит и на множестве $D$. По теореме Вейерштрасса непрерывная функция на непустом компактном множестве достигает наименьшего и наибольшего
    
    $\Rightarrow$ существует глобальный минимум и максимум.

    \item $L(x,y,\lambda) = x + y + \lambda h(x,y)$
    
    $L(x,y,\lambda) = x + y + \lambda (x^2 + y^2 - 2)$

    \subsection*{Система:}
    \[
    \begin{cases}
    \dfrac{\partial L}{\partial x} = 1 + 2\lambda x \\
    \dfrac{\partial L}{\partial y} = 1 + 2\lambda y \\
    \dfrac{\partial L}{\partial \lambda} = x^2 + y^2 - 2
    \end{cases}
    \qquad
    \begin{cases}
    1 + 2\lambda x = 0 \\
    1 + 2\lambda y = 0 \\
    x^2 + y^2 - 2 = 0
    \end{cases} \quad -
    \]
\end{enumerate}

\newpage

\[
(1 + 2\lambda x) - (1 + 2\lambda y) = 0
\]
\[
2\lambda x - 2\lambda y = 0
\]
\[
2\lambda(x - y) = 0
\]

\begin{align*}
2\lambda &= 0 & x - y &= 0 \\
\lambda &= 0 & x &= y \\
1 + 2 \cdot 0 \cdot x &= 0 & x^2 + y^2 &= 2 \\
1 &\neq 0 & x^2 + x^2 &= 2 \\
\Rightarrow \lambda &\neq 0 & 2x^2 &= 2 \\
& & x &= 1 \quad x = -1
\end{align*}

Т.к. $x = y$, то $(x_1, y_1) = (1, 1)$

$(x_2, y_2) = (-1, -1)$

Из первого уравнения $1 + 2\lambda x = 0 \quad \lambda = -\dfrac{1}{2x}$

$(1, 1) : \quad \lambda_1 = -\dfrac{1}{2 \cdot 1} = -\dfrac{1}{2}$

$(-1, -1) : \quad \lambda_2 = -\dfrac{1}{-2} = \dfrac{1}{2}$

$f(x,y) = x + y$

$f(1, 1) = 2 \qquad f(-1, -1) = -2$

$\Rightarrow$ \textbf{Наибольшее значение буде 2, наименьшее -2}

\end{document}
